% Options for packages loaded elsewhere
\PassOptionsToPackage{unicode}{hyperref}
\PassOptionsToPackage{hyphens}{url}
\PassOptionsToPackage{dvipsnames,svgnames,x11names}{xcolor}
\documentclass[
  11pt,
  a4paper,
]{article}
\usepackage{xcolor}
\usepackage[margin=19mm,headheight=15pt]{geometry}
\usepackage{amsmath,amssymb}
\setcounter{secnumdepth}{-\maxdimen} % remove section numbering
\usepackage{iftex}
\ifPDFTeX
  \usepackage[T1]{fontenc}
  \usepackage[utf8]{inputenc}
  \usepackage{textcomp} % provide euro and other symbols
\else % if luatex or xetex
  \usepackage{unicode-math} % this also loads fontspec
  \defaultfontfeatures{Scale=MatchLowercase}
  \defaultfontfeatures[\rmfamily]{Ligatures=TeX,Scale=1}
\fi
\usepackage{lmodern}
\ifPDFTeX\else
  % xetex/luatex font selection
\fi
% Use upquote if available, for straight quotes in verbatim environments
\IfFileExists{upquote.sty}{\usepackage{upquote}}{}
\IfFileExists{microtype.sty}{% use microtype if available
  \usepackage[]{microtype}
  \UseMicrotypeSet[protrusion]{basicmath} % disable protrusion for tt fonts
}{}
\makeatletter
\@ifundefined{KOMAClassName}{% if non-KOMA class
  \IfFileExists{parskip.sty}{%
    \usepackage{parskip}
  }{% else
    \setlength{\parindent}{0pt}
    \setlength{\parskip}{6pt plus 2pt minus 1pt}}
}{% if KOMA class
  \KOMAoptions{parskip=half}}
\makeatother
\setlength{\emergencystretch}{3em} % prevent overfull lines
\providecommand{\tightlist}{%
  \setlength{\itemsep}{0pt}\setlength{\parskip}{0pt}}

\usepackage{fontspec}
\setmainfont{Liberation Serif}
\setsansfont{Liberation Sans}
\setmonofont{Liberation Mono}[Scale=0.86]
\usepackage{ragged2e,indentfirst,xurl}
\usepackage{longtable,booktabs,array,calc,pdflscape}
\usepackage{xcolor,fancyhdr,titlesec,enumitem}
\definecolor{ink}{HTML}{19354A}
\definecolor{accent}{HTML}{167D8D}
\titleformat{\section}{\Large\sffamily\bfseries\color{ink}}{}{0pt}{}
\titleformat{\subsection}{\normalsize\sffamily\bfseries\color{ink}}{}{0pt}{}
\titleformat{\subsubsection}{\normalsize\sffamily\bfseries\color{ink}}{}{0pt}{}
\titlespacing*{\section}{0pt}{0pt}{8pt}
\titlespacing*{\subsection}{0pt}{8pt}{4pt}
\titlespacing*{\subsubsection}{0pt}{6pt}{3pt}
\setlength{\parindent}{1.25cm}
\setlength{\parskip}{3pt}
\setlength{\emergencystretch}{3em}
\setlength{\tabcolsep}{4pt}
\renewcommand{\arraystretch}{1.08}
\setlist{nosep,leftmargin=1.5em,topsep=3pt}
\pagestyle{fancy}\fancyhf{}
\fancyhead[L]{\small\sffamily ICSI405 / SEMINAR 02}
\fancyhead[R]{\small\sffamily Nogoolin / M2}
\fancyfoot[R]{\small\thepage}
\renewcommand{\headrulewidth}{0.3pt}
\AtBeginDocument{\fontsize{10.5}{12.5}\selectfont\justifying\setlength{\parindent}{1.25cm}}
\usepackage{bookmark}
\IfFileExists{xurl.sty}{\usepackage{xurl}}{} % add URL line breaks if available
\urlstyle{same}
\hypersetup{
  colorlinks=true,
  linkcolor={Maroon},
  filecolor={Maroon},
  citecolor={Blue},
  urlcolor={blue!55!black},
  pdfcreator={LaTeX via pandoc}}

\author{}
\date{}

\newcounter{none}
\begin{document}

\section{SRS Scope Charter}\label{doc1-srs-scope-charter}

\textbf{Nogoolin · ICSI405 / M2 / UE-2 · Тэнгис · 2026-09-15 · v0.1,
review draft}

\subsection{Төслийг сольсон
шийдвэр}\label{doc1-ux442ux4e9ux441ux43bux438ux439ux433-ux441ux43eux43bux44cux441ux43eux43d-ux448ux438ux439ux434ux432ux44dux440}

Би Seminar 1-ийн ажлаа зуны дадлагын Document RAG Chatbot төсөл дээр
хийж, хүлээлгэн өгсөн. Week 2-оос өөрийн бодит бүтээгдэхүүн, portfolio
болох Nogoolin-ийг үндсэн кейсээр сонгов. Энэ төсөлд шаардлага,
архитектур, API, deployment-ийн баримтыг одоогийн кодтой нийцүүлэх
хэрэгцээ байгаа тул M4, M5, M8, M9-ийн ажлыг бодит төсөлдөө үргэлжлүүлэх
боломжтой. Seminar 1-ийг дахин хийхгүй; шилжилтийн persona холбоог
M2-ийн нэмэлтээр ил тод тэмдэглэнэ.

\subsection{Зорилго ба
уншигч}\label{doc1-ux437ux43eux440ux438ux43bux433ux43e-ux431ux430-ux443ux43dux448ux438ux433ux447}

Каталогийн хэрэглэгч бүтээгдэхүүний тухай асуулга илгээх, хадгалсан
бүтээгдэхүүнээ дахин харах боломжийг тодорхойлно. Гол уншигч нь
хэрэгжүүлж, шалгах хөгжүүлэгч Тэнгис; хоёрдогч хичээлийн багш. Phase 0 баримтыг эх материал болгон ашиглаж,
технологийн сонголтыг SDD-д, шалгаж болох үр дүнг SRS-д тусгана.

\subsection{Included --- энэ
хүрээнд}\label{doc1-included-ux44dux43dux44d-ux445ux4afux440ux44dux44dux43dux434}

\begin{itemize}
\tightlist
\item
  \textbf{Inquiry + wishlist:} guest болон нэвтэрсэн хэрэглэгчийн
  асуулга, хүлээн авсан баталгаа; нэвтэрсэн хэрэглэгчийн
  хадгалах/харах/хасах урсгал. Админ inbox нь асуулга очсоныг шалгах
  төгсгөлийн цэг.
\item
  \textbf{Өгөгдлийн тусгаарлалт:} хэрэглэгч өөрийн wishlist, inquiry
  history-г л уншиж/өөрчилнө; админ эрхийг тусдаа шалгана.
\item
  \textbf{Hero fallback:} 3D asset бэлэн бус үед орлуулагч, WebGL-гүй
  болон хөдөлгөөн багасгах тохиргоонд статик хувилбар; каталогт хүрэх
  боломж.
\item
  \textbf{Шалгах ба build хийх хүрээ:} API-ийн local Docker build,
  deploy-оос өмнөх тестийн зайлшгүй хаалт. Docker-ийн сонголт SDD-д,
  амжилт/алдааны шалгуур NFR-03-д байна.
\end{itemize}

\subsection{Excluded --- M2/M3-ийн энэ багцаас
хассан}\label{doc1-excluded-m2m3-ux438ux439ux43d-ux44dux43dux44d-ux431ux430ux433ux446ux430ux430ux441-ux445ux430ux441ux441ux430ux43d}

\textbf{Бүтээгдэхүүн бүрийн 360° viewer:} тус бүрд нь 3D asset бэлтгэх
ажил ганцаараа хөгжүүлэх MVP-д хэт их тул энгийн зургийн gallery
ашиглана (D-07). Энэ нь нүүрний 3D intro-г хассан шийдвэр биш.
Viewer-ийг post-MVP-д дахин авч үзнэ; энэ багцад хэрэгжүүлэх, тестлэх
үүрэг үүсгэхгүй.

\subsection{Postponed --- дараа
хийх}\label{doc1-postponed-ux434ux430ux440ux430ux430-ux445ux438ux439ux445}

\textbf{Phase 5:} захиалга, төлбөр/checkout, хүргэлтийн ажиллагаа; одоо
\nolinkurl{delivery_enabled=false} гэсэн шийдвэр хүчинтэй. \textbf{Phase
6:} cloud Supabase холбох, production deployment баталгаажуулах, App
Store/Play Store нийтлэх. Mobile inquiry/wishlist, бодит GLB/.riv asset
болон бүх каталог/admin шаардлагыг энэ таван шаардлага бүрэн хамрахгүй;
бүтээгдэхүүний backlog-д үлдээнэ.

\subsection{Хүрээ тогтоох үндэслэл ба дуусах
нөхцөл}\label{doc1-ux445ux4afux440ux44dux44d-ux442ux43eux433ux442ux43eux43eux445-ux4afux43dux434ux44dux441ux43bux44dux43b-ux431ux430-ux434ux443ux443ux441ux430ux445-ux43dux4e9ux445ux446ux4e9ux43b}

Chinchilla-ийн х. 53--55 дахь аргаас: \textbf{хэн юу хүлээж байна?} ---
хэрэглэгч баталгаа, хөгжүүлэгч давтаж шалгах нөхцөл; \textbf{ямар далд
таамаг байна?} --- нэвтэрсэн л бол бүх өгөгдөл харах эрхтэй биш, asset
болон cloud бэлэн гэж үзэхгүй; \textbf{эхлээд юуг баримтжуулах вэ?} ---
эхлэлээс үр дүн хүртэлх inquiry/wishlist урсгал. Эдгээр нь номын аргыг
төсөлдөө хэрэглэсэн асуултууд, үгчилсэн эшлэл биш.

M2-ийн бэлэн багц: SRS/SDD загвар, 5 шаардлага, 8 баганат traceability,
persona холбоо, review draft. Ярилцлага хийгдээгүй; peer review илгээх,
нийтлэх/submission алхмыг хэрэглэгчийн заавраар алгассан. W1 persona-ийн
шинэ төсөлд шилжүүлсэн холбоог reviewer-ээр баталгаажуулах шаардлагатай.

\textbf{Эх:} \href{../../phase-0/01-vision.md}{Vision},
\href{../../phase-0/10-roadmap.md}{Roadmap},
\href{../../../agent-context/DECISIONS.md}{D-07 / ADR-011},
\href{Software_Project_Documentation_Week_02_Handout.pdf}{W2 handout};
Chinchilla, бүлэг 5, хэвлэмэл х. 53--55 (локал номын PDF х. 74--76).

\clearpage
\section{Nogoolin persona --- Week 2
нэмэлт}\label{doc2-nogoolin-persona-week-2-ux43dux44dux43cux44dux43bux442}

\textbf{Тэнгис · 2026-09-15 · US-2.3 · Өөрийн ажиглалт; ярилцлага
хийгдээгүй.}

\subsection{Уншигч ба бүтээгдэхүүний
хэрэглэгч}\label{doc2-ux443ux43dux448ux438ux433ux447-ux431ux430-ux431ux4afux442ux44dux44dux433ux434ux44dux445ux4afux4afux43dux438ux439-ux445ux44dux440ux44dux433ux43bux44dux433ux447}

\textbf{Хөгжүүлэгч persona:} Тэнгис, TypeScript/React ашигладаг,
web/API/өгөгдлийн санг холбож ажиллуулна. Зорилго нь Nogoolin-ийн
inquiry/wishlist урсгалыг local орчинд давтаж шалгаж, алдааг аль
давхаргаас хайхаа ойлгох. Хэрэгсэл: Fedora, VS Code, Git, pnpm, Docker,
DevTools.

\textbf{Бүтээгдэхүүний хэрэглэгч:} каталог үзэж, асуулга илгээх эсвэл
дараа үзэх бүтээгдэхүүнээ хадгалах хүн. Энэ нь
\href{../../phase-0/01-vision.md}{Phase 0 vision}-ийн catalog
visitor/inquiry customer тайлбарын хураангуй; шинэ хүнтэй ярилцсан үр
дүн биш.

\subsection{P-NG-01 --- Session ба өгөгдөл эзэмших эрхийг
ялгах}\label{doc2-p-ng-01-session-ux431ux430-ux4e9ux433ux4e9ux433ux434ux4e9ux43b-ux44dux437ux44dux43cux448ux438ux445-ux44dux440ux445ux438ux439ux433-ux44fux43bux433ux430ux445}

Нэвтрэлт ажиллаж байсан ч A хэрэглэгч B-ийн wishlist, inquiry history-г
харах ёсгүй. Шалгах зааварт guest, customer A/B, admin-ийн ялгаа байхгүй
бол эрхийн алдаа анзаарагдахгүй үлдэнэ. Энэ бол код болон security
баримтаас гаргасан эрсдэлийн ажиглалт; бодит халдлага болсон гэсэн үг
биш.

\textbf{Source:}
\href{../../phase-0/08-security.md\#9-idor-prevention}{08-security §9},
\href{../../../backend/api/src/repositories/supabase/cart.repository.ts}{cart
repository},
\href{../../../backend/api/src/repositories/supabase/inquiry.repository.ts}{inquiry
repository}. \textbf{Холбоо:}
\href{../../requirements/requirements.md\#nfr-01}{NFR-01}.

\subsection{P-NG-02 --- Asset бэлэн бус үед урсгалыг
шалгах}\label{doc2-p-ng-02-asset-ux431ux44dux43bux44dux43d-ux431ux443ux441-ux4afux435ux434-ux443ux440ux441ux433ux430ux43bux44bux433-ux448ux430ux43bux433ux430ux445}

Бодит Green Tara GLB/.riv бэлэн болоогүй үед каталогийн урсгалыг
хүлээлгэхгүй шалгах хэрэгтэй. Asset тохируулаагүй нөхцөл дэх орлуулагч,
WebGL-гүй болон reduced-motion хувилбарын үр дүнг тус бүр тодорхой
болгоно.

\textbf{Source:} \href{../../../agent-context/ASSET_SPECS.md}{Asset
contract},
\href{../../../apps/web/src/components/intro/deity.tsx}{deity.tsx},
\href{../../../apps/web/src/components/intro/hero-intro.tsx}{hero-intro.tsx}.
\textbf{Холбоо:}
\href{../../requirements/requirements.md\#nfr-02}{NFR-02}.

\subsection{P-NG-03 --- Хуучин заавар ба бодит орчин
зөрөх}\label{doc2-p-ng-03-ux445ux443ux443ux447ux438ux43d-ux437ux430ux430ux432ux430ux440-ux431ux430-ux431ux43eux434ux438ux442-ux43eux440ux447ux438ux43d-ux437ux4e9ux440ux4e9ux445}

Flutter/Express гэсэн хуучин тайлбаруудыг одоогийн React Native/Fastify
шийдвэртэй андуурч болно. Мөн local өгөгдлийн сан байхгүй үед тест
алгасагдсаныг бүрэн шалгалт өнгөрсөн гэж ойлгох эрсдэлтэй.

\textbf{Source:}
\href{../../../agent-context/DECISIONS.md}{ADR-003/004/011},
\href{../../../backend/api/tests/inquiry-cart.integration.test.ts}{integration
test-ийн skip салаа}. \textbf{Холбоо:}
\href{../../requirements/requirements.md\#nfr-03}{NFR-03},
\href{../../architecture/sdd-template.md}{SDD}.

\subsection{Week 1-тэй холбосон
үндэслэл}\label{doc2-week-1-ux442ux44dux439-ux445ux43eux43bux431ux43eux441ux43eux43d-ux4afux43dux434ux44dux441ux43bux44dux43b}

\href{../sem1/wiki/02_audience_persona.md}{Өмнө хүлээлгэн өгсөн W1
persona}-ийн \textbf{P1} нь нэвтрэлтийн cookie/session-ийг шалгах
зааврын хэрэгцээ байсан. Түүний authentication-ийг оношлох хэрэгцээг
Nogoolin-д өргөжүүлж \textbf{P1 → P-NG-01 → NFR-01} гэж холбов.
Authentication болон ownership нь өөр ойлголт; W1-д IDOR байсан гэж
үзээгүй. Мөн W1-ийн \textbf{P3} нь local ба pipeline-д өөр үр дүн
гарахыг оношлох, өөр замтай байх хэрэгцээ; үүнээс \textbf{P3 → P-NG-03 →
NFR-03} холбоог шинэ төсөлд шилжүүлэн гаргав.

Эдгээр нь Week 2-т нэмсэн тайлбарласан холбоо. US-2.3-ийн W1 pain point
шалгуурт энэ шилжүүлэлт нийцэх эсэхийг багш/reviewer-ээр баталгаажуулна;
Week 1-ийн эхийг өөрчлөөгүй.

\clearpage
\section{SRS Template --- Nogoolin}\label{doc3-srs-template-nogoolin}

\textbf{ICSI405 / M2 / US-2.1 · Тэнгис · 2026-09-15 · v0.1 / draft}

\subsection{1. Оршил}\label{doc3-ux43eux440ux448ux438ux43b}

\textbf{1.1 Purpose.} Энэ SRS нь inquiry, wishlist, хувийн өгөгдлийн
тусгаарлалт, hero fallback, API-ийн тестийн хаалтын хүлээгдэх үр дүнг
хөгжүүлэгч болон reviewer-т тодорхойлно. M2-д таван шаардлагын draft
бэлтгэсэн; бүрэн бүтээгдэхүүний SRS эсвэл ISO нийцлийн гэрчилгээ биш.

\textbf{1.2 Scope.} \href{../../requirements/scope-charter.md}{Scope
Charter}-ийн included/excluded/postponed шийдвэрийг мөрдөнө. Catalog нь
inquiry/wishlist-ийн урьдчилсан нөхцөл; 360° product viewer хасагдсан,
Phase 5/6 ажиллагаа хойшлогдсон. Үндсэн төсөл W2-оос Nogoolin болсон.

\textbf{1.3 Definitions.} Inquiry = бүтээгдэхүүний асуулга; wishlist =
дараа үзэхээр хадгалсан жагсаалт; guest = нэвтрээгүй хэрэглэгч; customer
= өөрийн мэдээлэлд эрхтэй хэрэглэгч; admin = тусгай удирдах эрхтэй
хэрэглэгч. IDOR = эзэмшигчийг шалгалгүй өөр хүний мэдээлэлд хүрэх алдаа;
RLS = өгөгдлийн мөрийн түвшний хандалтын дүрэм. FR/NFR болон
Test/Inspection нь \href{../../requirements/requirements.md}{шаардлагын
баримт}-д тайлбартай.

\textbf{1.4 References.}
\href{Software_Project_Documentation_Week_02_Handout.pdf}{W2 handout},
Lecture 02, слайд 4--8; \href{../../phase-0/README.md}{Phase 0 index},
\href{../../requirements/traceability-matrix.md}{кодын нотолгоо}.
\href{https://ieeexplore.ieee.org/document/720574}{IEEE 830-1998} нь
SRS-ийн бүтэц/чанарын зөвлөмж;
\href{https://www.iso.org/standard/72089.html}{ISO/IEC/IEEE 29148:2018}
нь requirements engineering-ийн процесс, мэдээллийн агуулгыг
тодорхойлдог. Доорх нь хичээлийн зориулалтаар тохируулсан IEEE 830
бүтэц; чанарын найман шалгуурыг Lecture 02-ын rubric-аар ашиглав.

\textbf{1.5 Overview.} §2 нь бүтээгдэхүүний орчин, хэрэглэгч, таамаг;
§3-ын mapping нь бүрэн/дутуу хэсэг;
\href{../../requirements/requirements.md}{5 шаардлага},
\href{../../requirements/traceability-matrix.md}{traceability},
\href{../../architecture/sdd-template.md}{SDD} нь дэлгэрэнгүй холбоотой
баримтууд.

\subsection{2. Ерөнхий
тодорхойлолт}\label{doc3-ux435ux440ux4e9ux43dux445ux438ux439-ux442ux43eux434ux43eux440ux445ux43eux439ux43bux43eux43bux442}

\textbf{2.1 Product perspective.} Nogoolin бол шашин, зан үйлийн
бүтээгдэхүүний каталог. Web хэрэглэгч бүтээгдэхүүн үзэж, асуулга
илгээнэ; customer хадгалсан бүтээгдэхүүн болон асуулгын түүхээ харна;
admin inbox-оос асуулгыг боловсруулна. Нэг хэрэглэгчийн хувийн
мэдээллийг бусдад үзүүлэхгүй.

\textbf{2.2 Product functions.} Энэ draft-ийн үйлдэл FR-01 (асуулга),
FR-02 (wishlist); тэдгээрийн чанарын хязгаар NFR-01\ldots03.
Бүтээгдэхүүн засварлах, search, auth-ийн бүх нарийвчилсан FR нь Phase 0
эхэд байгаа бөгөөд W3-д сонгон өргөжүүлнэ.

\textbf{2.3 User characteristics.}
\href{wiki/persona-and-pivot.md}{Persona нэмэлт}: хэрэглэгч энгийн
монгол интерфэйс, баталгаа хүснэ; хөгжүүлэгч давтаж шалгах алхам, эрхийн
ялгаа, орчны хамаарлыг мэдэх хэрэгтэй. Эдгээр нь баримт/кодын шинжилгээ;
хэрэглэгчийн ярилцлага pending.

\textbf{2.4 Constraints.} Ганцаараа хөгжүүлэх хүрээ; монгол интерфэйс;
delivery идэвхгүй; cloud холболт Phase 6 хүртэл хойшлогдсон.
Fastify/Next.js/Supabase/Docker-ийн сонголт нь
\href{../../architecture/sdd-template.md}{SDD}-д байна.

\textbf{2.5 Assumptions/dependencies.} Local өгөгдлийн сан, нийтлэгдсэн
бүтээгдэхүүн болон test customer/admin бэлэн үед acceptance шалгалт
ажиллана. Бодит GLB/.riv asset, cloud account, production log бэлэн гэж
үзэхгүй. Тестүүд ажилласан эсэхийг үр дүнгээс тусад нь шалгана.

\textbf{2.6 Apportioning.} W3-д өргөжүүлсэн SRS, W4-д архитектур, W5-д
API, W8-д диаграм, W9-д Docs-as-Code/CI evidence бэлтгэх төлөвлөгөөтэй.
Эдгээр нь хичээлийн долоо хоног; Phase 0--6 нь бүтээгдэхүүний roadmap
тул хооронд нь адилтгахгүй.

\subsection{3. IEEE 830 section
mapping}\label{doc3-ieee-830-section-mapping}

\nolinkurl{complete} = M2 загварт тухайн хэсгийн агуулга/холбоос байна;
хэрэгжилт болон stakeholder approval гэсэн үг биш.
\nolinkurl{planned (Wxx)} = нөхөх долоо хоног.
\nolinkurl{n-a (justified)} = энэ хүрээнд хамаарахгүй, шалтгаан нь
Content-д бий. Мөр бүр яг нэг төлөвтэй.

\begin{longtable}[]{@{}
  >{\raggedright\arraybackslash}p{(\linewidth - 4\tabcolsep) * \real{0.2600}}
  >{\raggedright\arraybackslash}p{(\linewidth - 4\tabcolsep) * \real{0.5500}}
  >{\raggedright\arraybackslash}p{(\linewidth - 4\tabcolsep) * \real{0.1900}}@{}}
\toprule\noalign{}
\begin{minipage}[b]{\linewidth}\raggedright
Section
\end{minipage} & \begin{minipage}[b]{\linewidth}\raggedright
Content / эх материал
\end{minipage} & \begin{minipage}[b]{\linewidth}\raggedright
Status
\end{minipage} \\
\midrule\noalign{}
\endhead
\bottomrule\noalign{}
\endlastfoot
1.1 Purpose & §1.1: хөгжүүлэгч/reviewer-ийн зорилго & complete \\
1.2 Scope & Scope Charter; 01 vision, 10 roadmap & complete \\
1.3 Definitions & §1.3: inquiry, wishlist, IDOR, RLS & complete \\
1.4 References & §1.4: handout, lecture, Phase 0, код & complete \\
1.5 Overview & §1.5: холбоотой баримтын бүтэц & complete \\
2.1 Product perspective & §2.1; 01 vision, 03 use cases & complete \\
2.2 Product functions & §2.2; FR-01/02-ийн хүрээ & complete \\
2.3 User characteristics & Persona нэмэлт; 01 vision, W1 persona &
complete \\
2.4 Constraints & §2.4; Scope Charter, SDD & complete \\
2.5 Assumptions/dependencies & §2.5; 10 roadmap, asset/DB нөхцөл &
complete \\
2.6 Apportioning & §2.6; Phase 5/6 postponed & complete \\
3.1.1 User interfaces & §3.1 доор; 07 wireframes; persona & complete \\
3.1.2 Hardware interfaces & Тусгай төхөөрөмж/сенсортой холболт энэ
web/API хүрээнд байхгүй. & n-a (justified) \\
3.1.3 Software interfaces & §3.1 доор; 06 API, shared schema; wishlist
OpenAPI gap & planned (W05) \\
3.1.4 Communication interfaces & §3.1: JSON, auth, status; 06 API, 08
security & complete \\
3.2 Functional requirements & FR-01/02 draft; M3-д сонгосон хүрээнд 15
FR болгон өргөжүүлэх & planned (W03) \\
3.3 Performance requirements & 02 requirements-ийн NFR-PER; reference
орчин/ачааллыг хэмжих нөхцөлтэй болгох & planned (W03) \\
3.4 Logical database requirements & 04 ER эх + inquiry/customer/cart
migration; өгөгдөл/unique/ownership дүрмийг нийцүүлэх & planned (W03) \\
3.5 Design constraints & §2.4 ба SDD; delivery-off, local-first, solo
хүрээ & complete \\
3.6.1 Reliability & NFR-02 draft; recovery, asset 404 шалгуурыг нөхөх &
planned (W03) \\
3.6.2 Availability & 02 requirements; outage хэмжих нөхцөл тодорхойлох &
planned (W03) \\
3.6.3 Security & NFR-01 draft; 08 security; 5 NFR багцад өргөжүүлэх &
planned (W03) \\
3.6.4 Maintainability & NFR-03 gap; 09 deployment; CI-ийн бодит run
evidence & planned (W09) \\
3.6.5 Portability & 02 requirements; browser/mobile дэмжлэгийн шалгах
матриц & planned (W03) \\
3.7 Other requirements & Audit/хувийн мэдээллийн хадгалах ба устгах
бодлого тодруулах; 08 security & planned (W03) \\
Appendix A: Traceability & 8 багана, 5 ID, source ба persona холбоо &
complete \\
Appendix B: Assumptions/quality & §2.5 ба §4; эх сурвалжийн хязгаарлалт
& complete \\
\end{longtable}

\subsubsection{3.1 External interface-ийн M2
эхлэл}\label{doc3-external-interface-ux438ux439ux43d-m2-ux44dux445ux43bux44dux43b}

\textbf{User:} бүтээгдэхүүний detail-ээс тусдаа inquiry page рүү орж
нэр/утас, сонголтот зурвас илгээнэ; амжилтын баталгаа ба 400/404/429
алдааг ялгана. Wishlist-д хадгалах/хасах, profile-д зөвхөн өөрийн
асуулга харах; guest хамгаалагдсан хуудсанд нэвтрэх урсгалтай. Hero нь
skip/static fallback-тай. Persona-ийн P-NG-01 нь эрхийн алдаа, P-NG-02
нь asset-гүй төлөвийг тайлбарлах хэрэгцээг өгнө.

\textbf{Software/communication:} REST JSON интерфэйс:
\nolinkurl{/api/v1/inquiries}, \nolinkurl{/api/v1/cart},
\nolinkurl{/api/v1/inquiries/mine}; protected API нь Bearer session
token-оор хэрэглэгчийг тогтооно. 201/204 амжилт, 400 validation, 401
нэвтрэлт, 403 эрх, 404 олдоогүй, 429 rate limit. Production transport нь
HTTPS гэсэн security эхийн хязгаартай; localhost нь хөгжүүлэлтийн орчин.
Field/schema-ийн бүрэн contract W5; 06 OpenAPI-д cart/history нэмэлт
хараахан нийцээгүй.

\subsection{4. Чанарын review ба W3 ажлын
хүрээ}\label{doc3-ux447ux430ux43dux430ux440ux44bux43d-review-ux431ux430-w3-ux430ux436ux43bux44bux43d-ux445ux4afux440ux44dux44d}

\begin{longtable}[]{@{}
  >{\raggedright\arraybackslash}p{(\linewidth - 2\tabcolsep) * \real{0.2500}}
  >{\raggedright\arraybackslash}p{(\linewidth - 2\tabcolsep) * \real{0.7500}}@{}}
\toprule\noalign{}
\begin{minipage}[b]{\linewidth}\raggedright
Lecture 02-ын шалгуур
\end{minipage} & \begin{minipage}[b]{\linewidth}\raggedright
Энэ draft-д хэрэглэсэн арга
\end{minipage} \\
\midrule\noalign{}
\endhead
\bottomrule\noalign{}
\endlastfoot
Correct & Source нь бодит файл; interview/approval pending гэж
тэмдэглэсэн \\
Unambiguous & Input, action, output, pass/fail-ийг тус бүр нэрлэсэн \\
Complete & Таван шаардлагын нөхцөл/хязгаар бий; бүтээгдэхүүний бүрэн SRS
W3 \\
Consistent & Email-гүй inquiry, wishlist ≠ checkout; технологи SDD-д \\
Ranked & Шаардлага бүр must; optional боломж энэ багцад ороогүй \\
Verifiable & T-01\ldots03, I-01\ldots02; нотолгооны төлөв тусдаа \\
Modifiable & Тогтвортой ID, салгасан Markdown, нэг traceability эх \\
Traceable & 8 багана; Phase 0/код/тайлбарласан W1 холбоо \\
\end{longtable}

\clearpage
\section{Таван шалгаж болох
шаардлага}\label{doc4-ux442ux430ux432ux430ux43d-ux448ux430ux43bux433ux430ux436-ux431ux43eux43bux43eux445-ux448ux430ux430ux440ux434ux43bux430ux433ux430}

\textbf{Nogoolin · M2 / US-2.2 · Тэнгис · 2026-09-15 · v0.1 / draft}

FR нь үйлдлийн, NFR нь чанар/хязгаарлалтын шаардлага. Доорх ID нь энэ
шинэ багцын ID; Phase 0-ийн урт ID-г солихгүй. \nolinkurl{must} =
зайлшгүй, \nolinkurl{shall} = тохирсон хүрээнд заавал биелүүлэх,
\nolinkurl{may} = сонголтот гэсэн хичээлийн тэмдэглэгээг хэрэглэв. Энд
тавуулаа \nolinkurl{must}. Owner: \textbf{Тэнгис}. Source interview:
\textbf{хийгдээгүй}; эх нь repo-ийн файл ба өөрийн шинжилгээ. Pass/fail
нь төлөвлөсөн шалгуур; тестийг энэ ажлаар ажиллуулаагүй.

\subsection{FR-01}\label{doc4-fr-01}

\textbf{Асуулга хадгалах · Priority: must · Verification: Test (T-01)}

Зочин эсвэл нэвтэрсэн хэрэглэгч зөв нэр, утас, сонголтот зурвас болон
бүтээгдэхүүний ID-тай асуулга илгээхэд систем нэг шинэ асуулга хадгалж,
\nolinkurl{201}, давтагдашгүй ID, \nolinkurl{status=new} бүхий баталгааг
\textbf{must} буцаана.

\textbf{T-01:} local өгөгдлийн санд нийтлэгдсэн бүтээгдэхүүн, rate
limit-д хүрээгүй IP бэлтгэнэ. Нэр \nolinkurl{Тэнгис}, тестийн утас
\nolinkurl{99112233}, сонголтот зурвас болон бүтээгдэхүүний ID-тай
\nolinkurl{POST /api/v1/inquiries} явуулна. Guest болон customer нөхцөл
бүрд 201, нэг мөр нэмэгдсэн, status=new, guest-ийн customer\_id=null,
customer-ийнх өөрийн ID байвал pass. Нэр 1 тэмдэгт эсвэл утас
\nolinkurl{abc} үед 400, мөр нэмэгдэхгүй; байхгүй бүтээгдэхүүний зөв
хэлбэртэй UUID үед 404, мөр нэмэгдэхгүй байна. Аль нэг хүлээлт зөрвөл
fail. Утасны бүрэн формат нь
\href{../../../packages/validation-schemas/src/common.ts}{shared
schema}-д байна.

\textbf{Source:} \href{../../phase-0/02-requirements.md}{Phase 0
FR-INQ-001\ldots003},
\href{../../../packages/validation-schemas/src/inquiry.schema.ts}{inquiry
schema},
\href{../../../backend/api/src/controllers/inquiry.controller.ts}{controller}.
Email талбар оруулахгүй; message API дээр сонголтот. \textbf{Төлөв:}
хэрэгжилт, тестийн эх байгаа; шинэ ажиллуулалтын нотолгоо pending.

\subsection{FR-02}\label{doc4-fr-02}

\textbf{Wishlist-ийн төлөв хадгалах · Priority: must · Verification:
Test (T-02)}

Нэвтэрсэн хэрэглэгч нийтлэгдсэн бүтээгдэхүүнийг хадгалах, жагсаалтаа
харах, хасахад систем хэрэглэгч-бүтээгдэхүүн хос бүрд хамгийн ихдээ нэг
хадгалсан бичлэгтэй төлөвийг \textbf{must} хадгална.

\textbf{T-02:} A хэрэглэгч бүтээгдэхүүн P-г
\nolinkurl{POST /api/v1/cart}-аар хоёр удаа нэмэхэд хүсэлт бүр 201,
\nolinkurl{GET /api/v1/cart}-д P яг нэг удаа байна. Хуудсыг дахин
ачаалахад P хадгалагдсан байна. \nolinkurl{DELETE /api/v1/cart/P}-г хоёр
удаа дуудахад хоёул 204, дараах GET-д P байхгүй байвал pass. Guest
хүсэлт 401, draft/байхгүй бүтээгдэхүүн нэмэх хүсэлт 404 байна. Аль нэг
хүлээлт зөрвөл fail. Энэ нь wishlist бөгөөд тоо хэмжээ, нийт үнэ,
checkout агуулахгүй.

\textbf{Source:}
\href{../../../backend/api/src/controllers/cart.controller.ts}{cart
controller},
\href{../../../backend/api/src/repositories/supabase/cart.repository.ts}{cart
repository}, \href{../../../apps/web/src/app/wishlist/page.tsx}{wishlist
page},
\href{../../../supabase/migrations/20260726120000_inquiry_customer_and_cart.sql}{20260726120000
migration}. \textbf{Төлөв:} хэрэгжилт, тестийн эх байгаа; runtime
pending. Phase 0-д wishlist-ийн бие даасан FR байгаагүй; одоогийн кодоос
сэргээн тодорхойлов.

\subsection{NFR-01}\label{doc4-nfr-01}

\textbf{Хэрэглэгчдийн өгөгдлийг тусгаарлах · Priority: must ·
Verification: Inspection (I-01)}

Customer A-ийн session-ээр wishlist болон inquiry history
унших/өөрчлөхөд систем өөр customer B-ийн хувийн мөрийг хариунд
оруулахгүй, өөрчлөхгүй байхыг \textbf{must} хангана; guest нь эдгээр
хамгаалагдсан үйлдэлд 401 авна. Админ inbox-д зөвхөн баталгаажсан admin
хүрнэ.

\textbf{I-01:} хамгаалагдсан cart GET/POST/DELETE, inquiry mine GET,
admin inquiry GET/PATCH-ийн controller → service → repository болон
RLS-ийг шалгана. Эзэмшигчийн ID баталгаажсан session-ээс ирдэг,
хэрэглэгчийн body/query-гээс ирдэггүй; cart нь \nolinkurl{user_id},
inquiry нь \nolinkurl{customer_id}-аар шүүгддэг; insert нь session-ийн
ID оноодог; RLS болон admin role guard бий гэдгийг мөр бүрд нотолбол
pass. Нэг хамгаалагдсан зам ownership/role шалгалтгүй бол fail. Public
catalog, guest inquiry create-д сохроор \nolinkurl{user_id} filter
шаардахгүй.

\textbf{Дэмжих тест:} A/B хоёр customer үүсгээд A-ийн хариунд B-ийн ID
огт байхгүй, A хасахад B-ийн хадгалсан бүтээгдэхүүн хэвээр, customer
admin inbox-д 403 авдгийг шалгана. \textbf{Source:}
\href{../../phase-0/08-security.md\#9-idor-prevention}{security §9},
\href{wiki/persona-and-pivot.md}{P-NG-01 / W1 холбоо}, дээрх хоёр
repository ба migration. \textbf{Төлөв:} статик замыг шалгасан; бүрэн
I-01 checklist, хоёр-customer runtime нотолгоо pending.

\clearpage
\subsection{NFR-02}\label{doc4-nfr-02}

\textbf{Asset-гүй үед каталогт хүрэх · Priority: must · Verification:
Test (T-03)}

Нүүрний 3D asset тохируулаагүй үед систем орлуулагчийг, WebGL боломжгүй
эсвэл reduced-motion идэвхтэй үед статик hero-г \textbf{must} үзүүлж,
хэрэглэгчийг каталогт хүрэх боломжтой байлгана.

\textbf{T-03:} desktop 1440×900, mobile хэмжээ 390×844 бүхий browser-д
sessionStorage-г цэвэрлэж дараахыг тус бүр шалгана: (a) GLB URL
тохируулаагүй --- орлуулагч харагдах; (b) WebGL unavailable --- статик
hero; (c) reduced-motion --- хөдөлгөөнт intro алгасагдах. Intro гарсан
нөхцөлд skip нь navigation эхэлснээс 1 секундийн дотор харагдаж, дарахад
home төлөвт орно. Нөхцөл бүрд бүтээгдэхүүний холбоосоор
\nolinkurl{/products} нээгдэж, жагсаалт харагдвал pass; хоосон дэлгэц,
түгжигдсэн scroll, хүрэх боломжгүй каталог байвал fail. Browser/version,
local build, хэмжсэн хугацааг тэмдэглэнэ. Буруу URL/404 asset-ийн
recovery нь тусдаа gap, энэ шаардлагын батлагдсан хэрэгжилт гэж үзэхгүй.

\textbf{Source:}
\href{../../phase-0/07-uiux-wireframes.md}{WF-INTRO-03/05/06/07},
\href{../../../apps/web/src/components/intro/hero-intro.tsx}{hero-intro},
\href{../../../apps/web/src/components/intro/deity.tsx}{deity},
\href{wiki/persona-and-pivot.md}{P-NG-02}. \textbf{Төлөв:} код байгаа;
шинэ browser evidence pending.

\subsection{NFR-03}\label{doc4-nfr-03}

\textbf{Deploy-оос өмнөх тестийн хаалт · Priority: must · Verification:
Inspection (I-02)}

API production deploy нь тухайн commit-ийн шаардлагатай integration
тестүүд бодитоор ажиллаж амжилттай дууссан, production container build
амжилттай болсон үед л эхлэхийг систем \textbf{must} хангана. Тест fail
болох, бүх integration тест алгасагдах, эсвэл Docker build fail болох
нөхцөл бүр deploy-г хаана.

\textbf{I-02:} CI тохиргоо, тестийн entry point, job dependency болон
гурван сөрөг тохиолдлын CI log-ийг шалгана. Test fail / DB unavailable
буюу all-skipped / Docker build fail бүрд deploy job skipped эсвэл
blocked, амжилттай тохиолдолд deploy нь зөвхөн хоёр шалгалтын дараа
eligible болсон байвал pass. Энэ review нь production deploy ажиллуулах
шаардлагагүй. Auto-deploy зэрэг CI-ийг тойрох зам Phase 6-д хаагдсан
эсэхийг мөн шалгана.

\textbf{Source:}
\href{../../phase-0/02-requirements.md}{NFR-MAIN-005/006},
\href{../../phase-0/09-deployment.md}{deployment §8.4},
\href{../../../.github/workflows/api-deploy.yml}{api workflow},
\href{../../../backend/api/tests/inquiry-cart.integration.test.ts}{test
skip branch}, \href{../../../backend/api/Dockerfile}{Dockerfile}.
\textbf{Төлөв: gap / planned (W09).}
\nolinkurl{lint-and-test → docker-build → deploy} холбоо байгаа ч
тестүүд DB-гүй үед skip хийнэ, workflow local Supabase асаахгүй. Иймд
хаалт бүрэн хэрэгжсэн гэж үзэхгүй. Web workflow-д test команд байхгүй;
энэ NFR API-д хамаарна, web рүү өргөтгөх нь дараагийн шаардлага.

\clearpage
\begin{landscape}
\section{Traceability Matrix}\label{doc5-traceability-matrix}

\textbf{Nogoolin · M2 / US-2.3 · v0.1 · 2026-09-15}

Доорх нь яг 8 баганатай, 5 мөртэй. Last-Reviewed нь баримт/кодыг уншсан
огноо; тест ажиллуулсан огноо биш. Owner нь Тэнгис (хөгжүүлэгч, баримт
бичигч). Бүх шаардлага draft, stakeholder interview хийгдээгүй.

\begin{longtable}[]{@{}
  >{\raggedright\arraybackslash}p{(\linewidth - 14\tabcolsep) * \real{0.0650}}
  >{\raggedright\arraybackslash}p{(\linewidth - 14\tabcolsep) * \real{0.2550}}
  >{\raggedright\arraybackslash}p{(\linewidth - 14\tabcolsep) * \real{0.0650}}
  >{\raggedright\arraybackslash}p{(\linewidth - 14\tabcolsep) * \real{0.1200}}
  >{\raggedright\arraybackslash}p{(\linewidth - 14\tabcolsep) * \real{0.1150}}
  >{\raggedright\arraybackslash}p{(\linewidth - 14\tabcolsep) * \real{0.1450}}
  >{\raggedright\arraybackslash}p{(\linewidth - 14\tabcolsep) * \real{0.1400}}
  >{\raggedright\arraybackslash}p{(\linewidth - 14\tabcolsep) * \real{0.0950}}@{}}
\toprule\noalign{}
\begin{minipage}[b]{\linewidth}\raggedright
ID
\end{minipage} & \begin{minipage}[b]{\linewidth}\raggedright
Source
\end{minipage} & \begin{minipage}[b]{\linewidth}\raggedright
Owner
\end{minipage} & \begin{minipage}[b]{\linewidth}\raggedright
Verification
\end{minipage} & \begin{minipage}[b]{\linewidth}\raggedright
Dependency
\end{minipage} & \begin{minipage}[b]{\linewidth}\raggedright
Risk
\end{minipage} & \begin{minipage}[b]{\linewidth}\raggedright
Status
\end{minipage} & \begin{minipage}[b]{\linewidth}\raggedright
Last-Reviewed
\end{minipage} \\
\midrule\noalign{}
\endhead
\bottomrule\noalign{}
\endlastfoot
\href{../../requirements/requirements.md\#fr-01}{FR-01} &
\href{../../phase-0/02-requirements.md}{FR-INQ-001\ldots003};
\href{../../../packages/validation-schemas/src/inquiry.schema.ts}{schema};
\href{../../../backend/api/src/controllers/inquiry.controller.ts}{controller}
& Тэнгис & Test: T-01;
\href{../../../backend/api/tests/inquiry-cart.integration.test.ts}{тестийн
эх} & Бүтээгдэхүүн, local DB, validation, rate limit & Email/message-ийн
хуучин баримтын зөрүү; давхардсан илгээлт & Draft; runtime pending &
2026-09-15 \\
\href{../../requirements/requirements.md\#fr-02}{FR-02} &
\href{../../../backend/api/src/controllers/cart.controller.ts}{cart
controller};
\href{../../../supabase/migrations/20260726120000_inquiry_customer_and_cart.sql}{migration}
& Тэнгис & Test: T-02;
\href{../../../backend/api/tests/inquiry-cart.integration.test.ts}{тестийн
эх} & Auth, published product, unique pair, NFR-01 & Wishlist-ийг
checkout-тай андуурах & Draft; runtime pending & 2026-09-15 \\
\href{../../requirements/requirements.md\#nfr-01}{NFR-01} &
\href{../sem1/wiki/02_audience_persona.md}{W1 persona P1} →
\href{wiki/persona-and-pivot.md\#p-ng-01--session-ба-өгөгдөл-эзэмших-эрхийг-ялгах}{P-NG-01};
\href{../../phase-0/08-security.md\#9-idor-prevention}{08 §9};
\href{../../../backend/api/src/repositories/supabase/cart.repository.ts}{cart
repo};
\href{../../../backend/api/src/repositories/supabase/inquiry.repository.ts}{inquiry
repo} & Тэнгис & Inspection: I-01; A/B runtime нэмэлт & Session
identity, ownership filter, RLS, RBAC & Өөр хүний мэдээлэл алдагдах; W1
холбооны зөвшөөрөл pending & Draft; partial inspection & 2026-09-15 \\
\href{../../requirements/requirements.md\#nfr-02}{NFR-02} &
\href{wiki/persona-and-pivot.md\#p-ng-02--asset-бэлэн-бус-үед-урсгалыг-шалгах}{P-NG-02};
\href{../../phase-0/07-uiux-wireframes.md}{wireframes};
\href{../../../apps/web/src/components/intro/hero-intro.tsx}{hero} &
Тэнгис & Test: T-03; browser evidence pending & Hero, fallback, catalog
route & Asset байхгүйгээс хэрэглэгч каталогт хүрэхгүй & Draft; runtime
pending & 2026-09-15 \\
\href{../../requirements/requirements.md\#nfr-03}{NFR-03} &
\href{../sem1/wiki/02_audience_persona.md}{W1 P3} →
\href{wiki/persona-and-pivot.md\#p-ng-03--хуучин-заавар-ба-бодит-орчин-зөрөх}{P-NG-03};
\href{../../phase-0/09-deployment.md}{09 deployment};
\href{../../../.github/workflows/api-deploy.yml}{workflow};
\href{../../../backend/api/tests/inquiry-cart.integration.test.ts}{skip
branch} & Тэнгис & Inspection: I-02; CI negative-case logs pending &
Local DB in CI, test exit status, Docker build & Бүх тест skip боловч
deploy eligible болох & Gap; planned (W09) & 2026-09-15 \\
\end{longtable}

\subsection{Эх сурвалжийн үнэн зөв
байдал}\label{doc5-ux44dux445-ux441ux443ux440ux432ux430ux43bux436ux438ux439ux43d-ux4afux43dux44dux43d-ux437ux4e9ux432-ux431ux430ux439ux434ux430ux43b}

W1-ийн P1 нь cookie/session асуудал бөгөөд IDOR биш. NFR-01 рүү холбосон
нь Week 2-т нэмсэн, тайлбарласан шинэ хэрэглээ. W1 P3-ийн local/pipeline
ялгаанаас NFR-03-ийн тестийн орчны эрсдэлийг гаргасан. Багш эдгээр
холбоог зөвшөөрсөн гэж үзээгүй;
\href{wiki/peer-review-draft.md}{peer-review draft}-д шалгуулах асуулт
бий. Ярилцлагын source зохиогоогүй.

\subsection{Шинэчлэх
дүрэм}\label{doc5-ux448ux438ux43dux44dux447ux43bux44dux445-ux434ux4afux440ux44dux43c}

Шаардлага өөрчлөгдвөл ID-г хадгалж, эх холбоос, verification-ийн
pass/fail, dependency, risk, status, огноог хамтад нь шинэчилнэ. Тестийн
run URL/log нэмсний дараа л verification status-ийг verified болгоно.
Confluence-ийн Requirements хэсэгт энэ файлын GitHub permalink-ийг
оруулахад бэлэн; хэрэглэгчийн заавраар нийтлэлт/submission-ийг алгассан,
remote commit хийгээгүй.

\end{landscape}
\clearpage
\section{SDD Template --- Nogoolin}\label{doc6-sdd-template-nogoolin}

\textbf{ICSI405 / M2 · Тэнгис · 2026-09-15 · v0.1 / draft}

\subsection{1. Зорилго ба шаардлагын
холбоо}\label{doc6-ux437ux43eux440ux438ux43bux433ux43e-ux431ux430-ux448ux430ux430ux440ux434ux43bux430ux433ux44bux43d-ux445ux43eux43bux431ux43eux43e}

Энэ SDD нь \href{../../requirements/requirements.md}{таван шаардлагыг}
хэрэгжүүлэх бүтэц, технологийг тайлбарлана. M2-д одоогийн кодоос үндсэн
бүтэц, design decision, gap-ийг тогтоов; arc42-ийн дэлгэрэнгүй W4,
OpenAPI W5, диаграм W8-д төлөвлөгдсөн. SRS нь хэрэглэгчийн авах үр дүн,
SDD нь түүнийг хэрэгжүүлэх арга.

\subsection{2. Системийн орчин ба
бүрэлдэхүүн}\label{doc6-ux441ux438ux441ux442ux435ux43cux438ux439ux43d-ux43eux440ux447ux438ux43d-ux431ux430-ux431ux4afux440ux44dux43bux434ux44dux445ux4afux4afux43d}

\begin{longtable}[]{@{}
  >{\raggedright\arraybackslash}p{(\linewidth - 4\tabcolsep) * \real{0.2600}}
  >{\raggedright\arraybackslash}p{(\linewidth - 4\tabcolsep) * \real{0.5500}}
  >{\raggedright\arraybackslash}p{(\linewidth - 4\tabcolsep) * \real{0.1900}}@{}}
\toprule\noalign{}
\begin{minipage}[b]{\linewidth}\raggedright
Бүрэлдэхүүн
\end{minipage} & \begin{minipage}[b]{\linewidth}\raggedright
Шийдвэр / бодит байршил
\end{minipage} & \begin{minipage}[b]{\linewidth}\raggedright
Холбогдох шаардлага
\end{minipage} \\
\midrule\noalign{}
\endhead
\bottomrule\noalign{}
\endlastfoot
Web + admin & Next.js/TypeScript; \nolinkurl{apps/web}; inquiry,
wishlist, profile, admin inbox & FR-01/02, NFR-01/02 \\
API & Fastify/TypeScript; \nolinkurl{backend/api}; \nolinkurl{/api/v1},
local port 3001 & FR-01/02, NFR-01 \\
Өгөгдөл/Auth & Supabase local PostgreSQL/Auth/Storage;
\nolinkurl{supabase/migrations} & FR-01/02, NFR-01 \\
Shared validation & Zod; \nolinkurl{packages/validation-schemas} &
FR-01/02 \\
Hero & Three.js/R3F; \nolinkurl{apps/web/src/components/intro} &
NFR-02 \\
Build/CI & pnpm workspaces, Docker,
\nolinkurl{.github/workflows/api-deploy.yml} & NFR-03 \\
\end{longtable}

\textbf{Шийдвэрийн эх:}
\href{../../../agent-context/DECISIONS.md}{ADR-003/004/005/006/007/008/011}.
Mobile нь React Native + Expo prebuild/dev client; энэ M2 багц mobile
inquiry/wishlist-ийг дууссан гэж үзэхгүй. Railway/Vercel/cloud Supabase
нь Phase 6-ын байршуулалтын зорилт, одоо ажиллаж буй production орчин
гэж батлаагүй.

\subsection{3. Backend-ийн дотоод
бүтэц}\label{doc6-backend-ux438ux439ux43d-ux434ux43eux442ux43eux43eux434-ux431ux4afux442ux44dux446}

Controller → Service → Repository. Controller нь request/response ба
validation hook; service нь бизнесийн дүрэм; repository нь өгөгдөлд
хүрэх кодыг хариуцна. Auth hook хэрэглэгчийн session-ийг тогтооно.
Бизнесийн логикийг route руу, DB query-г service рүү холихгүй.

\textbf{FR-01:} inquiry controller → inquiry service → inquiry
repository → \nolinkurl{inquiries}. \textbf{FR-02:} cart controller →
cart service → cart repository → \nolinkurl{cart_items}. Shared schema-г
web/API нэг эхээс хэрэглэнэ; client validation нь server validation-ийг
орлохгүй.

\subsection{4. Өгөгдөл ба эрхийн
дизайн}\label{doc6-ux4e9ux433ux4e9ux433ux434ux4e9ux43b-ux431ux430-ux44dux440ux445ux438ux439ux43d-ux434ux438ux437ux430ux439ux43d}

\nolinkurl{inquiries.customer_id} нь нэвтэрсэн илгээгчийг заана, guest-д
null; шинэ асуулгын төлөв \nolinkurl{new}. \nolinkurl{cart_items} нь
\nolinkurl{(user_id, product_id)} unique хостой. Нэмэх үйлдлийг давтахад
нэг бичлэг хэвээр; хасах нь давтаж дуудаж болох үйлдэл. Эх:
\href{../../../supabase/migrations/20260726120000_inquiry_customer_and_cart.sql}{20260726120000
migration},
\href{../../../backend/api/src/repositories/supabase/cart.repository.ts}{cart
repository},
\href{../../../backend/api/src/repositories/supabase/inquiry.repository.ts}{inquiry
repository}.

\textbf{NFR-01:} customer ID-г request body-оос авахгүй, баталгаажсан
session-ээс дамжуулна. Cart read/delete нь
\nolinkurl{.eq('user_id', userId)}, inquiry mine нь
\nolinkurl{.eq('customer_id', customerId)}; insert session-ийн ID
онооно. RLS нь DB түвшний хамгаалалт, admin scope нь RBAC шалгалттай.
Бүх query-д нэг ижил \nolinkurl{user_id} шүүлтүүр шаардахгүй: public
catalog, guest create, admin ажиллагаа тус тусын эрхийн дүрэмтэй.

\clearpage
\subsection{5. Гадаад интерфэйс ба contract
gap}\label{doc6-ux433ux430ux434ux430ux430ux434-ux438ux43dux442ux435ux440ux444ux44dux439ux441-ux431ux430-contract-gap}

\href{../../phase-0/06-api-spec.yaml}{Phase 0 OpenAPI} нь эх contract.
Код дахь \nolinkurl{/cart} болон \nolinkurl{/inquiries/mine}, customer
ownership болон admin filter нэмэлтүүдийг W5-д тулган шинэчилнэ. Inquiry
нь email field-гүй; UI wireframe-ийн хуучин email мөрийг хэрэгжилттэй
ижил гэж үзэхгүй. \nolinkurl{INQ-YYYY-NNNN} нь UI-ийн derived display
утга бөгөөд дараалсан захиалгын дугаар биш.

\subsection{6. Hero fallback-ийн
дизайн}\label{doc6-hero-fallback-ux438ux439ux43d-ux434ux438ux437ux430ux439ux43d}

GLB URL тохируулаагүй үед \nolinkurl{deity.tsx} геометр орлуулагч
сонгоно. \nolinkurl{hero-intro.tsx} нь WebGL/reduced-motion-оор static
home төлөв сонгож, skip/catalog холбоосыг DOM дээр харуулна. Бодит
asset-ийг \href{../../../agent-context/ASSET_SPECS.md}{ASSET\_SPECS}-ийн
дагуу дараа солино. Алдаатай GLB URL-ийн error boundary/recovery-г энэ
draft баталгаажаагүй gap гэж тэмдэглэв. Product 360° viewer нь D-07-оор
post-MVP; нүүрний intro-оос тусдаа.

\subsection{7. Build, deploy ба
verification}\label{doc6-build-deploy-ux431ux430-verification}

API Dockerfile нь multi-stage build, non-root runtime ашигладаг.
Одоогийн CI dependency нь
\nolinkurl{lint-and-test → docker-build → deploy}.
\href{../../../backend/api/tests/inquiry-cart.integration.test.ts}{Тестийн
эх} DB хүрэхгүй үед skip хийдэг; workflow local Supabase бэлтгэхгүй тул
NFR-03-ын all-skipped хаалт хангагдаагүй. W9-д CI DB setup,
skip-fails-CI нөхцөл, гурван сөрөг шалгалтын log төлөвлөсөн. Production
auto-deploy нь энэ хаалтыг тойрохгүй байх тохиргоог Phase 6-д шалгана.

Энэ task-аар код, DB migration, CI-ийн ажиллагааг өөрчлөөгүй; шинэ
runtime/build тест ажиллуулаагүй.
\href{../../../agent-context/PROGRESS.md}{PROGRESS}-ийн өмнөх тестийн
тайлан нь түүхэн тэмдэглэл, одоогийн ажиллуулалтын нотолгоо биш.

\subsection{8. Үргэлжлүүлэн бөглөх
хэсгүүд}\label{doc6-ux4afux440ux433ux44dux43bux436ux43bux4afux4afux43bux44dux43d-ux431ux4e9ux433ux43bux4e9ux445-ux445ux44dux441ux433ux4afux4afux434}

\begin{longtable}[]{@{}
  >{\raggedright\arraybackslash}p{(\linewidth - 4\tabcolsep) * \real{0.2600}}
  >{\raggedright\arraybackslash}p{(\linewidth - 4\tabcolsep) * \real{0.5500}}
  >{\raggedright\arraybackslash}p{(\linewidth - 4\tabcolsep) * \real{0.1900}}@{}}
\toprule\noalign{}
\begin{minipage}[b]{\linewidth}\raggedright
Хэсэг
\end{minipage} & \begin{minipage}[b]{\linewidth}\raggedright
Нөхөх агуулга
\end{minipage} & \begin{minipage}[b]{\linewidth}\raggedright
Status
\end{minipage} \\
\midrule\noalign{}
\endhead
\bottomrule\noalign{}
\endlastfoot
Architecture views & arc42 context, building blocks, runtime,
deployment; FR/NFR холбоо & planned (W04) \\
API contracts & cart/history schema ба бодит error/status нийцүүлэлт &
planned (W05) \\
Diagrams & C4, ER, inquiry/auth sequence source ба review & planned
(W08) \\
Design verification & CI negative-case evidence, Docker build evidence,
traceability шинэчлэлт & planned (W09) \\
\end{longtable}

\textbf{Эх:} \href{../../phase-0/README.md}{Phase 0 index},
\href{../../../agent-context/DECISIONS.md}{шийдвэрүүд},
\href{../../requirements/srs-template.md}{SRS},
\href{../../requirements/traceability-matrix.md}{traceability}.

\clearpage
\section{Peer review draft --- M2}\label{doc7-peer-review-draft-m2}

\textbf{Тэнгис · Nogoolin · 2026-09-15 · Бэлтгэсэн; илгээгээгүй}

\subsection{Reviewer-т зориулсан
тайлбар}\label{doc7-reviewer-ux442-ux437ux43eux440ux438ux443ux43bux441ux430ux43d-ux442ux430ux439ux43bux431ux430ux440}

Week 2-оос Document RAG Chatbot-оос Nogoolin руу шилжсэн. Week 1-ийн
ажлыг өөрчлөөгүй. Phase 0 баримтаас SRS-ийн шаардлага болон SDD-ийн
технологийн шийдвэрийг салгаж, таван шаардлагын draft ба traceability
бэлтгэлээ.

\textbf{Унших дараалал:}
\href{../../requirements/scope-charter.md}{Scope Charter} →
\href{wiki/persona-and-pivot.md}{persona нэмэлт} →
\href{../../requirements/srs-template.md}{SRS загвар} →
\href{../../requirements/requirements.md}{5 шаардлага} →
\href{../../requirements/traceability-matrix.md}{traceability} →
\href{../../architecture/sdd-template.md}{SDD}.

\subsection{Шалгах
асуултууд}\label{doc7-ux448ux430ux43bux433ux430ux445-ux430ux441ux443ux443ux43bux442ux443ux443ux434}

\begin{enumerate}
\def\labelenumi{\arabic{enumi}.}
\tightlist
\item
  Included/excluded/postponed нь ганцаараа хөгжүүлэх хүрээнд хангалттай
  тодорхой байна уу? Product 360° viewer ба нүүрний hero хоёр ялгагдаж
  байна уу?
\item
  Mapping-ийн мөр бүр яг нэг төлөвтэй, n-a нь тайлбартай байна уу?
  Complete төлөвийг template-ийн агуулгатай гэж зөв ойлгож байна уу?
\item
  FR-01/02, NFR-01\ldots03 бүрд оролт, үйлдэл, гаралт, pass/fail хязгаар
  хангалттай юу? Олон шалгуурыг M3-д дэд шаардлага болгох хэрэгтэй юу?
\item
  \textbf{W1 P1 → P-NG-01 → NFR-01}, \textbf{W1 P3 → P-NG-03 → NFR-03}
  гэсэн тайлбарласан холбоо нь төсөл сольсон нөхцөлд US-2.3-ын persona
  шаардлагыг хангах уу?
\item
  Ownership-ийг session-ээс авах дүрэм ойлгомжтой юу? Guest inquiry,
  customer history, admin inbox-ийг зөв ялгасан уу?
\item
  NFR-03-т all-skipped тестийг fail гэж үзэх шалгуур хангалттай юу?
  SRS/SDD-ийн хил зааг алдагдсан өгүүлбэр байна уу?
\end{enumerate}

\subsection{Санал бүртгэх
загвар}\label{doc7-ux441ux430ux43dux430ux43b-ux431ux4afux440ux442ux433ux44dux445-ux437ux430ux433ux432ux430ux440}

\begin{longtable}[]{@{}
  >{\raggedright\arraybackslash}p{(\linewidth - 6\tabcolsep) * \real{0.2000}}
  >{\raggedright\arraybackslash}p{(\linewidth - 6\tabcolsep) * \real{0.1900}}
  >{\raggedright\arraybackslash}p{(\linewidth - 6\tabcolsep) * \real{0.3500}}
  >{\raggedright\arraybackslash}p{(\linewidth - 6\tabcolsep) * \real{0.2600}}@{}}
\toprule\noalign{}
\begin{minipage}[b]{\linewidth}\raggedright
Reviewer / огноо
\end{minipage} & \begin{minipage}[b]{\linewidth}\raggedright
Requirement/section
\end{minipage} & \begin{minipage}[b]{\linewidth}\raggedright
Асуудал ба санал
\end{minipage} & \begin{minipage}[b]{\linewidth}\raggedright
Хариу / өөрчлөлт
\end{minipage} \\
\midrule\noalign{}
\endhead
\bottomrule\noalign{}
\endlastfoot
Pending & Pending & Санал хараахан ирээгүй & Pending \\
\end{longtable}

\subsection{Source interview ба
submission}\label{doc7-source-interview-ux431ux430-submission}

Source interview хийгдээгүй; owner нь Тэнгис, source нь repo-ийн файл ба
өөрийн ажиглалт. Хүлээн авагч болон Confluence/LMS холбоос байхгүй тул
2026-09-15-ны хэрэглэгчийн заавраар илгээлт/нийтлэлийг алгассан. Peer
review хийсэн, neighboring team-д илгээсэн гэж тэмдэглэхгүй. Дараа нь
хүлээн авах суваг тодорхой болбол Confluence Requirements хэсэгт repo
permalink холбож, submission evidence хадгалж болно.

\subsection{Богино
retrospective}\label{doc7-ux431ux43eux433ux438ux43dux43e-retrospective}

\textbf{Сайн болсон:} Phase 0-ийн их хэмжээний эхээс M2-д хэрэгтэй жижиг
хүрээ сонгож, кодтой тулгав. \textbf{Тодорхойгүй байсан:} SRS-ийн ``юу
хийх'' ба SDD-ийн ``яаж хийх'' өгүүлбэрүүд холилдсон; W1 persona нь өөр
төсөл байсан. \textbf{Дараагийн өөрчлөлт:} шаардлага бүрийг source →
observable result → verification дарааллаар бичиж, тестийн команд
амжилттай дууссан эсэхээс гадна тест үнэхээр ажилласан эсэхийг шалгана.

Chinchilla-ийн бүлэг 5-аас далд таамаг ба одоогийн дуусах нөхцөлийг
эхэлж тодруулах аргыг үргэлжлүүлэн хэрэглэнэ. Номын үйл явц ба энэ ажлын
хамгийн том зөрүү нь бодит хэрэглэгч/reviewer-ийн feedback хараахан
аваагүй байхад draft-ыг кодын шинжилгээнээс эхлүүлсэн явдал юм.


\end{document}
